% !PRINTED 151
\weberpage{151}
%<<CONT
the coefficients were independent variables; but they cannot
become false when definite numerical values are put in their
place, since, at least for the integral symmetric functions, no
denominators occur that could vanish.

We first make some applications to particular symmetric
functions which are important for what follows. Already in
\S.\,19 it was remarked that the difference product

\weq{1}{\begin{aligned}
P = {}&(\alpha_1 - \alpha_2)\,(\alpha_1 - \alpha_3)\ldots(\alpha_1 - \alpha_n)\\
&\phantom{(\alpha_1 - \alpha_2)\,}(\alpha_2 - \alpha_3)\ldots(\alpha_2 - \alpha_n)\\
&\cdot\;\;\cdot\;\;\cdot\;\;\cdot\;\;\cdot\;\;\cdot\;\;\cdot\;\;\cdot\;\;\cdot\;\;\cdot\;\;\cdot\;\;\cdot\\
&\phantom{(\alpha_1 - \alpha_2)\,(\alpha_1 - \alpha_3)\ldots}(\alpha_{n-1} - \alpha_n)
\end{aligned}}

merely changes its sign under an interchange of two elements.
The square of this product will therefore remain unchanged
under all permutations of two of the variables $\alpha$, and hence
also under all permutations; that is, it will be a symmetric function
of the variables. The order of this symmetric function
is, as one sees from (1), equal to

\[
(2\,n - 2,\ \ 2\,n - 4 \ldots 2,\, 0).
\]

If, then, we assume the function $f(x)$, whose roots are the $\alpha_1, \alpha_2 \ldots \alpha_n$,
in the form

\[
f(x) = a_0\,x^n + a_1\,x^{n-1} + a_2\,x^{n-2} + \cdots + a_{n-1}\,x + a_n,
\]

so that the highest power of $x$ does not necessarily have the coefficient
1, then $P^2$ is an integral rational function of

\weq{2}{\frac{a_1}{a_0},\ \ \frac{a_2}{a_0} \ldots \frac{a_n}{a_0},}

which, on multiplication by $a_0^{2n-2}$, goes over into a \gesperrt{homogeneous} function
of $a_0, a_1, \ldots a_n$. We call this homogeneous function,
which is thus defined by

\weq{3}{a_0^{2\,n-2}\,P^2 = D}

the \gesperrt{discriminant of the function} $f(x)$. If we set
$a_0 = 1$, we obtain the \gesperrt{inhomogeneous form of the
discriminant}\footnote{The discriminant is, apart from the factor $(-\,1)^{\frac{n\,(n-1)}{2}}\ a_0^{2\,n-2}$,
the same as what \gesperrt{Gauss} calls the determinant of the function (l. c., art. 6).}.
